% Included by main.tex. Requires amsmath, amssymb, amsthm and theorem,
% proposition, lemma, remark environments. All notation is defined locally.
\section{Prediction movement, gain, and the evidence needed to grow}
\label{sec:movement-theory}

The information used here is evidence about a proposed change in task loss.
It is not the entropy of a neuron. Two candidates can have the same small
positive gain and very different sampling difficulty, because they move
predictions by different amounts. We first make this distinction precise,
then exhibit a whole sequence for which discarding evidence between
structural decisions imposes a quadratic label cost. The latter cost is a
property of a specified audit interface; we also prove how it can be avoided.
The decomposition and statistical ingredients are classical. The explicit module construction illustrates the consequences of retaining
or discarding off-action observations; its statistical mechanism is classical.

\subsection{The geometry of a paired Brier comparison}

Let $Y\in\{0,1\}$ and $\eta(x)=\mathbb P(Y=1\mid X=x)$.
Predictions $p,q:\mathcal X\to[0,1]$ are fixed before independent audit data.
Write
\begin{equation}
 D=(Y-p(X))^2-(Y-q(X))^2,\quad h=q-p,\quad
 \Delta=\mathbb ED,\quad S=\mathbb E h^2.
 \label{eq:movement-def}
\end{equation}
Positive $\Delta$ means that $q$ improves on $p$; $S$ measures prediction
movement on population inputs.

\begin{proposition}[Movement and paired variance]
\label{prop:movement}
One has
\begin{align}
 D&=h(2Y-p-q), & \Delta&=2\mathbb E[h(\eta-p)]-S,\label{eq:gain-movement}\\
 \operatorname{Var}(D)
 &=4\mathbb E[h^2\eta(1-\eta)]
   +\operatorname{Var}\{h(2\eta-p-q)\}.\label{eq:variance-movement}
\end{align}
In particular, $|D|\le1$ and $\operatorname{Var}(D)\le4S$.
If $a\le\eta\le1-a$ almost surely, for $0<a\le1/2$, then
$\operatorname{Var}(D)\ge4a(1-a)S$.
If $p=\mathbb E[Y\mid\mathcal A]$ and
$q=\mathbb E[Y\mid\mathcal B]$ for sigma fields
$\mathcal A\subset\mathcal B$, then $\Delta=S$.
\end{proposition}
\begin{proof}
Expansion of the two squared losses gives \eqref{eq:gain-movement}.
Conditional on $X$, $h,p,q$ are constants, so
$\operatorname{Var}(D\mid X)=4h^2\eta(1-\eta)$ and
$\mathbb E[D\mid X]=h(2\eta-p-q)$.
The conditional-variance identity proves \eqref{eq:variance-movement}.
Both squared losses lie in $[0,1]$, giving $|D|\le1$; also
$|2Y-p-q|\le2$, hence $\operatorname{Var}(D)\le\mathbb ED^2\le4S$.
The noise lower bound follows from $\eta(1-\eta)\ge a(1-a)$.
Finally, $q-p$ is $\mathcal B$-measurable and
$\mathbb E[Y-q\mid\mathcal B]=0$, so
$\mathbb E[(Y-q)(q-p)]=0$. Substituting
$Y-p=(Y-q)+(q-p)$ in the difference of squared losses gives $\Delta=S$.
\end{proof}

The final assertion is the familiar projection identity for nested Bayes
refits. A trained nonlinear branch need not satisfy it. Its gain can be
small because two sizable terms in \eqref{eq:gain-movement} nearly cancel.
The ratio $S/\Delta$, when the population gain is positive, records this
discrepancy. It is a diagnostic, not a universal control objective.

\begin{proposition}[A movement-dependent sufficient audit size]
\label{prop:movement-upper}
Suppose an upper bound $S\le\bar S$ is fixed independently of the audit.
For $n$ iid paired observations and $u>0$, define
\[
 r_n=\sqrt{\frac{8\bar S u}{n}}+\frac{4u}{3n}.
\]
Each tail probability
$\mathbb P(\bar D_n-\Delta>r_n)$ and
$\mathbb P(\Delta-\bar D_n>r_n)$ is at most $e^{-u}$.
The test $\bar D_n>r_n$ therefore has false acceptance probability at most
$e^{-u}$ for $\Delta\le0$. When $\Delta>0$, the condition
\begin{equation}
 n>\max\left\{\frac{128\bar S u}{\Delta^2},
                  \frac{16u}{3\Delta}\right\}
 \label{eq:power-sample}
\end{equation}
gives power at least $1-e^{-u}$.
\end{proposition}
\begin{proof}
Put $Z=D-\Delta$ and $v=\mathbb EZ^2$. Then $|Z|\le2$.
If $v=0$, $D$ is almost surely constant and both tail bounds hold directly.
Assume henceforth $v>0$.
For $0\le\lambda<3/2$, expanding the exponential and using
$\mathbb E|Z|^k\le2^{k-2}v$ and
$k!\ge2\,3^{k-2}$ for $k\ge2$ yields
\[
 \log\mathbb Ee^{\lambda Z}
 \le\frac{\lambda^2v}{2(1-2\lambda/3)}.
\]
Apply Markov's inequality to the sum of $n$ independent copies and take
$\lambda=r/(v+2r/3)$ to obtain
\[
 \mathbb P(\bar Z_n\ge r)
 \le\exp\left\{-\frac{nr^2}{2(v+2r/3)}\right\}.
\]
For $r=\sqrt{2vu/n}+4u/(3n)$ the exponent has magnitude at least $u$.
Use $v\le4\bar S$ and repeat for $-Z$. Under
\eqref{eq:power-sample}, the square-root and linear terms of $r_n$ are
each smaller than $\Delta/4$, so $2r_n<\Delta$. Except on the lower-tail
failure event, $\bar D_n\ge\Delta-r_n>r_n$.
\end{proof}

This is a standard Bernstein argument with an explicit movement bound.
Replacing $\bar S$ by its same-sample empirical value without a confidence
bound would not satisfy the proposition. The empirical gate instead uses
an established empirical Bernstein inequality \cite{maurer2009}.
For nested Bayes refits $S=\Delta$, so the leading sufficient sample order
is $1/\Delta$; for unaligned changes it is $S/\Delta^2$.

\subsection{A matching pairwise lower-bound witness}

Write
$\operatorname{kl}(r,s)=r\log(r/s)+(1-r)\log\{(1-r)/(1-s)\}$.
For $\alpha,\beta\in(0,1)$ with $\alpha+\beta<1$, let
$b_{\alpha,\beta}=\operatorname{kl}(1-\beta,\alpha)$.

\begin{lemma}[Sequential testing, established change of measure]
\label{lem:stopped-kl}
Let $P,Q$ be finite observation laws with positive probabilities.
A randomized procedure observes iid draws, stops at a common stopping rule
$N$ that is almost surely finite under both laws, and outputs an event $A$
determined by its stopped transcript.
There is no initial information whose distribution differs between the laws.
If $P(A)\ge1-\beta$ and $Q(A)\le\alpha$, then
\[
 \mathbb E_PN\,\operatorname{KL}(P,Q)\ge b_{\alpha,\beta}.
\]
\end{lemma}
\begin{proof}
If $\mathbb E_PN=\infty$ there is nothing to prove. Otherwise the stopped
log-likelihood ratio is
$L_N=\sum_{i\ge1}1_{\{N\ge i\}}\log\{P(Z_i)/Q(Z_i)\}$.
The increments are bounded, and the sum of their absolute expectations
is finite because $\mathbb E_PN<\infty$. The indicator is determined
before draw $i$, whose log-likelihood increment has mean
$\operatorname{KL}(P,Q)$. Thus
$\mathbb E_PL_N=\mathbb E_PN\operatorname{KL}(P,Q)$.
The common stopping and randomization rules contribute no extra factor
to the likelihood ratio of the stopped transcript. Its KL divergence is
therefore $\mathbb E_PL_N$. Data processing to $1_A$ bounds this below by
$\operatorname{kl}(P(A),Q(A))$.
For $r>s$, binary KL increases with $r$ and decreases with $s$, as follows
by differentiating its displayed formula. Hence this last divergence is
at least $b_{\alpha,\beta}$.
\end{proof}
The lemma is the one-observation-channel case of the established
change-of-measure result of Kaufmann, Capp\'e, and Garivier
\cite[Lemma 1 and Appendix A.1]{kaufmann2016}; its proof is included to
make the information restriction inspectable.

\begin{theorem}[A movement-dependent necessary audit size]
\label{thm:movement-lower}
For every $0<d\le1/2$ and $0<\Delta\le d/2$, there are two constant
predictions with movement $S=d^2$ and two Bernoulli label laws such that
the proposed change is risk-neutral under the first law and improves
Brier risk by $\Delta$ under the second. Any audit with false acceptance
at most $\alpha$ under the first and power at least $1-\beta$ under the
second, with no initial distinguishing information, satisfies
\[
 \mathbb E_1N\ge b_{\alpha,\beta}\frac{S}{\Delta^2}.
\]
\end{theorem}
\begin{proof}
Use a one-point input space and predictions $p=(1-d)/2$, $q=(1+d)/2$.
Let $P_0$ have label mean $1/2$ and $P_1$ mean
$1/2+\Delta/(2d)\le3/4$. Then $D=d(2Y-1)$, so the gains are zero
and $\Delta$, while $S=d^2$.
The inequality $\log u\le u-1$ applied to both Bernoulli KL terms gives
$\operatorname{kl}(r,s)\le(r-s)^2/[s(1-s)]$.
Consequently $\operatorname{KL}(P_1,P_0)\le\Delta^2/d^2$.
Lemma~\ref{lem:stopped-kl} proves the claim.
\end{proof}

This is a worst-case pair construction with fixed predictions. It does not
assert that every neural pair with the same movement and gain has this
difficulty. In particular, the gain alone gives no universal sample bound:
some restricted noiseless comparisons have deterministic positive loss
differences. Also, a gate that always rejects is safe without any samples;
the lower bound depends essentially on the stated power requirement.

\section{The cost of resetting evidence across structural decisions}
\label{sec:reset-theory}

We now give a stationary neural sequence for which the event cost is
sharp and all gains add to a fixed risk budget. This is an audit theorem
with supplied modules. It is separate from the experiments that must learn
their own candidate directions and retrain all parameters.

\subsection{Population, neural actions, and observation interface}

Fix an even $K\ge2$, $0<\varepsilon\le1/4$, and known alternating signs
$s_t=(-1)^t$. Let $X$ be uniform on $\{1,\ldots,K\}$ and
\begin{equation}
 \mathbb P(Y=1\mid X=t)=\tfrac12+s_t\theta_t,
 \qquad 0\le\theta_t\le\varepsilon.
 \label{eq:module-law}
\end{equation}
With $\rho(z)=\max\{z,0\}$, define a three-ReLU module
\begin{equation}
 \phi_t(x)=\rho(x-t+1)-2\rho(x-t)+\rho(x-t-1).
 \label{eq:relu-module}
\end{equation}
At integer inputs $\phi_t(j)=1_{\{j=t\}}$.
An installed set $A$ predicts
$f_A=1/2+\varepsilon\sum_{t\in A}s_t\phi_t$, within $[1/4,3/4]$ on the
population support. The initial predictor is $f_\varnothing=1/2$.
Each action adds three actual hidden ReLUs. Their knots and amplitude
are supplied; this construction gives no guarantee that SGD finds them.

At opportunity $t$, module $t$ is proposed, whether or not earlier modules
were accepted. The controller may retain the existing network. Its paired
loss observation is
\begin{equation}
 D_t=1_{\{X=t\}}\bigl(2s_t\varepsilon(Y-1/2)-\varepsilon^2\bigr),
 \qquad g_t=\mathbb ED_t
   =\frac{2\varepsilon\theta_t-\varepsilon^2}{K}.
 \label{eq:module-difference}
\end{equation}
These expressions are independent of earlier acceptance decisions.
The known signs are part of the supplied candidate specification. When all
$\theta_t=\varepsilon$, the Bayes prediction alternates across inputs, so
an intercept alone cannot remove the excess risk. No claim of minimal
neural width is made: the theorem fixes the local-module action family.

The \emph{fresh paired-audit interface} supplies a new iid population stream
at each opportunity and reveals only the current $D_t$.
The controller may retain all previous paired observations, decisions,
and sampling times. It does not retain labels at inputs on which the
current action makes no prediction change, nor otherwise acquire data
outside this interface. Each audit may stop sequentially using the allowed
history and environment-independent randomness. Its sample count $N_t$
includes every labelled population example drawn, including zero paired
differences. Assume the following guarantees conditional on every possible
earlier allowed transcript:
\begin{enumerate}
 \item false acceptance is at most $\alpha$ whenever $g_t\le0$;
 \item acceptance probability is at least $1-\beta$ if
 $\theta_t=\varepsilon$;
 \item the stopping rule is almost surely finite under the tested laws,
 with finite expectation under the all-positive law (otherwise the
 lower bound is immediate).
\end{enumerate}
Here $\alpha,\beta\in(0,1)$ and $\alpha+\beta<1$.
These are distribution-uniform audit requirements, not assertions about
the power of the implemented neural gate.

\begin{theorem}[Serial evidence cost and a matching audit]
\label{thm:reset-cost}
Under \eqref{eq:module-law}, take $\theta_t=\varepsilon$ for every $t$.
Every candidate has gain and movement
$\Delta_t=S_t=\varepsilon^2/K$. The initial excess Brier risk above the
Bayes predictor is $D_0=\varepsilon^2$.
Every fresh paired-audit controller satisfying the preceding requirements
obeys
\begin{equation}
 \mathbb E\sum_{t=1}^K N_t
 \ge\frac{K^2}{2\varepsilon^2}b_{\alpha,\beta}
 =\frac{K^2}{2D_0}b_{\alpha,\beta}.
 \label{eq:quadratic-reset}
\end{equation}
The expected number accepted is at least $K(1-\beta)$, and expected risk
reduction is at least $(1-\beta)D_0$.
Conversely, the same interface admits an audit satisfying these guarantees
with
\begin{equation}
 \mathbb E\sum_{t=1}^K N_t
 =K^2\left\lceil\frac8{\varepsilon^2}
       \log\frac1{\min\{\alpha,\beta\}}\right\rceil.
 \label{eq:quadratic-upper}
\end{equation}
\end{theorem}
\begin{proof}
The identity
$\mathbb E(Y-f)^2=\mathbb E[\eta(1-\eta)]+\mathbb E(\eta-f)^2$
shows that each module removes exactly $\varepsilon^2/K$ excess risk.
All $K$ remove $D_0=\varepsilon^2$; prediction movement is also
$\varepsilon^2/K$.

Fix $t$ and compare the all-positive environment $P$ to $Q_t$, which only
replaces $\theta_t$ by $\varepsilon/2$. The stage-$t$ move is risk-neutral
under $Q_t$. Every earlier observation $D_s$, $s<t$, has exactly the same
law under both environments: it is zero at $X=t$, while its only
label-dependent value occurs at $X=s$, whose label law did not change.
Inducting through the common stopping and decision rules shows that the
entire earlier allowed transcript has the same law. Conditional on almost
every transcript under this common law, no initial information
distinguishes $P$ from $Q_t$.

The signed label $\widetilde Y_t=\{1+s_t(2Y-1)\}/2$, conditional on
$X=t$, has mean $1/2+\theta_t$. Relabelling by this invertible map
preserves KL. The one-sample full-pair divergence therefore satisfies
\begin{align*}
 \operatorname{KL}(P,Q_t)
 &=\frac1K\operatorname{kl}(1/2+\varepsilon,1/2+\varepsilon/2)\\
 &\le\frac1K\frac{(\varepsilon/2)^2}
                    {(1/2+\varepsilon/2)(1/2-\varepsilon/2)}
 \le\frac{2\varepsilon^2}{K}.
\end{align*}
The last denominator is at least $15/64$; the displayed factor two is
conservative. Mapping full pairs to $D_t$ cannot increase KL.
Apply Lemma~\ref{lem:stopped-kl} conditionally on the earlier transcript
to obtain $\mathbb E_P[N_t\mid\text{history}]
\ge Kb_{\alpha,\beta}/(2\varepsilon^2)$.
Integration and summation prove \eqref{eq:quadratic-reset}.
Conditional power and linearity of expectation give the remaining
lower-bound assertions. The decisions need not be independent.

For the upper bound, let
$m=\lceil8\varepsilon^{-2}\log(1/\min\{\alpha,\beta\})\rceil$.
At stage $t$, wait for $m$ nonzero paired observations; these are exactly
the examples with $X=t$, and either nonzero value identifies $Y$.
Accept if their mean signed label $\widetilde Y_t$ exceeds
$1/2+3\varepsilon/4$.
For $\theta_t\le\varepsilon/2$ and for $\theta_t=\varepsilon$, the mean
is at least $\varepsilon/4$ from the threshold in the appropriate
direction. The Bernoulli Hoeffding bound gives either error at most
$e^{-m\varepsilon^2/8}\le\min\{\alpha,\beta\}$.
For completeness, Hoeffding's bound follows by bounding the centered
Bernoulli moment generating function by $e^{\lambda^2/8}$,
applying Markov's inequality to $m$ independent samples, and optimizing
at $\lambda=4a$ for a deviation $a$; the result is $e^{-2ma^2}$.
The waiting time to $m$ occurrences of an event of probability $1/K$ has
mean $Km$, being the sum of $m$ geometric waiting times of mean $K$.
Summing over opportunities proves \eqref{eq:quadratic-upper}.
\end{proof}

For fixed small error probabilities, the cost is quadratic in the number
of useful proposed modules at fixed total removable risk. If all $K$
decisions must be simultaneously correct with probability at least
$1-\delta$, taking $\alpha=\beta=\delta/K$ supplies that guarantee and
adds the usual logarithmic factor. The lower bound is a change-of-measure
argument over an entire stationary neural sequence, not a repeated union
bound. Its limitation is equally central: the observation interface
discards off-action information, and the modules are prescribed.

\subsection{A constructive escape through shared evidence}

\begin{theorem}[Precommitted reuse removes a factor $K$]
\label{thm:shared-evidence}
In the same population, fix all modules before seeing labels and assume
each $\theta_t$ is either at most $\varepsilon/2$ or exactly $\varepsilon$.
One common batch of
\begin{equation}
 n=\left\lceil\frac{16K}{\varepsilon^2}\log\frac{2K}{\delta}\right\rceil
 \label{eq:shared-n}
\end{equation}
iid labelled pairs is sufficient to accept every positive separated
alternative and reject every nonpositive proposed change with probability
at least $1-\delta$. The accepted modules can be installed successively.
On the subfamily $\theta_t\in\{0,\varepsilon\}$ the resulting network
contains exactly the needed modules on this event, including no growth
when all $\theta_t=0$.
\end{theorem}
\begin{proof}
For each coordinate, use its signed labels $\widetilde Y_t$ in the common
batch and accept if their mean exceeds $1/2+3\varepsilon/4$; reject when
the batch contains no observations at that coordinate.
Let $C_t$ count examples at coordinate
$t$ and put $m_0=n/(2K)$. A multiplicative binomial bound gives
$\mathbb P(C_t<m_0)\le e^{-n/(8K)}$.
It follows directly from Markov's inequality applied to
$e^{-\lambda C_t}$: the binomial generating function is at most
$\exp\{(n/K)(e^{-\lambda}-1)\}$; taking $\lambda=\log2$ gives
an exponent no greater than $-n/(8K)$ at $C_t=n/(2K)$.
Conditional on counts and $C_t\ge m_0$, Hoeffding's bound makes the
test error at most
$e^{-m_0\varepsilon^2/8}\le\delta/(2K)$.
Also $e^{-n/(8K)}\le\delta/(2K)$ by \eqref{eq:shared-n} and
$\varepsilon\le1/4$. A union bound over the $K$ count failures and $K$
test errors proves the claim. Locality makes each comparison independent
of which other modules were installed, so sequential installation does
not change the statements being tested.
\end{proof}

The reuse theorem is valid because the local candidates and tests were
fixed independently of the labels. It does not validate arbitrary
holdout reuse after neural fitting or feedback-dependent proposal changes;
adaptive evaluation is a separate problem with established solutions and
limitations \cite{blum2015}. Alternatively, screening inputs without
requesting labels and labelling only $X=t$ reduces the separate tests to
$Km$ total labels. It still takes $K^2m$ expected input inspections unless
direct conditional input sampling is available. These resources should
not be conflated.

The exact sequence also explains why a rejected operation is not an
adequacy certificate. At $\theta_t=\varepsilon/2$, adding amplitude
$\varepsilon$ is risk-neutral, while adding amplitude $\varepsilon/2$
would improve risk. Our no-growth statement concerns a specified action
family; it certifies complete adequacy only on the subfamily
$\theta_t\in\{0,\varepsilon\}$ stated in
Theorem~\ref{thm:shared-evidence}.

\subsection{What transfers to a learned controller}

The elementary implication
\[
 \sum_{t=1}^K\Delta_t\le D_0,\quad
 \mathbb EN_t\ge c/\Delta_t
 \quad\Longrightarrow\quad
 \sum_{t=1}^K\mathbb EN_t\ge cK^2/D_0
\]
follows from Cauchy--Schwarz. The per-event premise is not automatic;
Theorem~\ref{thm:reset-cost} supplies an exact realization. For more general
comparisons, the relevant terms may be $S_t/\Delta_t^2$, and their observed
values should be reported instead of asserting a universal exponent.

There is an additional information-accounting caveat. If initial history
$H$ has different laws under the competing environments, the correct
testing constraint is
\begin{equation}
 \operatorname{KL}(P_H,Q_H)+
 \mathbb E_PN\operatorname{KL}(P,Q)
 \ge\operatorname{kl}(P(A),Q(A)),
 \label{eq:history-kl}
\end{equation}
when the new iid stream is independent of $H$ under each environment.
This follows from the KL chain rule and the proof of
Lemma~\ref{lem:stopped-kl}; history-dependent sampling replaces the second
term by a sum of conditional divergences. Initial training or retained
audit data can pay for the decision. Dropping the first term while allowing
arbitrary historical information would make a fresh-audit lower bound
false. The same labels cannot then be counted repeatedly as new purchases.

Consequently the neural experiments below test movement, realized gain,
and audit cost along ordinary trained trajectories; they are not numerical
proofs that those trajectories satisfy the local-module lower bound.
The observation distinction survives this limitation: failure to accept
can mean either an ineffective operation or inadequate evidence for a useful
one. Recording both terms is necessary to tell these explanations apart.
